\section{The Qubit}
The fundamental, indivisible unit of quantum information is the \textit{qubit} (quantum bit) \cite{nielsen2010quantum}.

While a classical bit is purely a logical abstraction realized physically by a macroscopic system with two distinct stable states (such as a high or low electrical voltage), a qubit is a microscopic, two-level quantum mechanical system.
Physical realizations of a qubit include the spin configuration of a single electron or the polarization state of a photon.

\subsection{Computational Basis}
Mathematically, the state of a single, isolated qubit is represented by a unit vector residing in a two-dimensional complex Hilbert space, $\mathbb{C}^2$.

We define an orthonormal basis $\{\ket{0}, \ket{1}\}$, designated as the \textit{computational basis}.
These basis vectors correspond directly to the classical, deterministic logical states $0$ and $1$:
\begin{equation}
    \ket{0} = \begin{pmatrix} 1 \\ 0 \end{pmatrix}, \quad \ket{1} = \begin{pmatrix} 0 \\ 1 \end{pmatrix}
\end{equation}

\subsection{Superposition}
Unlike a classical bit, which must be strictly in one state or the other, a qubit can exist in a linear combination of its basis states.
This is a fundamental property known as \textit{superposition}, which conceptually allows quantum algorithms to process vast amounts of information simultaneously.
A generalized qubit state $\ket{\psi}$ is mathematically described by the equation:
\begin{equation}
    \ket{\psi} = \alpha\ket{0} + \beta\ket{1}
\end{equation}
where $\alpha$ and $\beta$ are continuous complex numbers known as probability amplitudes.

According to the Born rule, if we actively measure this qubit in the computational basis, the delicate superposition will instantaneously collapse into one of the definite basis states.
The probability of observing the outcome $\ket{0}$ is exactly $P(0) = |\alpha|^2$, and the probability of observing $\ket{1}$ is exactly $P(1) = |\beta|^2$.
Because the qubit must physically exist in some state after measurement, these probabilities must sum to exactly 100\%.
This forces the state to be bound by the absolute normalization condition:
\begin{equation}
    |\alpha|^2 + |\beta|^2 = 1
\end{equation}

\subsection{The Bloch Sphere}
Because $\alpha$ and $\beta$ are complex numbers, a qubit state theoretically possesses four real degrees of freedom.
However, the normalization condition mathematically removes one of these degrees.
Furthermore, in quantum mechanics, the overall \textit{global phase} of a quantum state (a shared factor of $e^{i\gamma}$) has no observable physical consequence.
This is because all measurement probabilities are derived solely from the absolute squares of the amplitudes, effectively canceling out any global phase factor.
This physical reality removes a second degree of freedom.

Consequently, we can dramatically simplify the mathematical parameterization of any valid pure qubit state to rely on exactly two real angles, $\theta$ and $\phi$:
\begin{equation}
    \ket{\psi} = \cos\left(\frac{\theta}{2}\right)\ket{0} + e^{i\phi}\sin\left(\frac{\theta}{2}\right)\ket{1}
\end{equation}
Here, the polar angle $0 \leq \theta \leq \pi$ strictly dictates the probability distribution between the classical $\ket{0}$ and $\ket{1}$ outcomes.
The azimuthal angle $0 \leq \phi < 2\pi$ defines the \textit{relative phase} between the two states, a critically important quantum resource that drives wave interference in algorithms.

This elegant parameterization maps every possible pure single-qubit state to a unique geometric point on the surface of a three-dimensional unit sphere, known in quantum physics as the \textbf{Bloch Sphere} \cite{nielsen2010quantum}.
\begin{itemize}
    \item The North Pole ($\theta = 0$) geometrically represents the classical deterministic state $\ket{0}$.
    \item The South Pole ($\theta = \pi$) represents the classical deterministic state $\ket{1}$.
    \item Points lying perfectly on the equator (e.g., $\theta = \pi/2, \phi = 0$) represent perfectly equal superpositions.
    In these equatorial states, the qubit has exactly a 50\% chance of collapsing into either basis state upon measurement.
\end{itemize}

While the Bloch sphere provides an absolutely indispensable visual intuition for single-qubit operations—which mathematically correspond simply to 3D rotations on the surface of the sphere—it has a profound limitation.
It is critical to note that this helpful visualization completely fails to generalize to multi-qubit systems due to the complex, non-local nature of quantum entanglement.

\begin{figure}[htbp]
    \centering
    \includegraphics[width=0.5\textwidth]{Bloch_sphere.jpg}
    \caption{The Bloch sphere provides a geometric representation of the pure state space of a single qubit. Source: Wikipedia.}
    \label{fig:Bloch_sphere}
\end{figure}
